(1) 滑块速度为 \((\dot X+\dot s\cos\theta,-\dot s\sin\theta)\)，所以
\[
\mathcal L=\frac{M+m}{2}\dot X^2+m\cos\theta\,\dot X\dot s+\frac m2\dot s^2
-\frac k2s^2+mg s\sin\theta .
\]
\(X\) 未显含于 \(\mathcal L\)，故水平动量
\[
P_X=(M+m)\dot X+m\cos\theta\,\dot s
\]
恒定；依初态有 \(P_X=0\)，因此 \(\dot X=-m\cos\theta\,\dot s/(M+m)\)。

(2) 消去 \(\dot X\) 后动能为 \(T=\frac12m_{\rm eff}\dot s^2\)，其中
\[
m_{\rm eff}=m\left(1-\frac{m\cos^2\theta}{M+m}\right)
=\frac{m(M+m\sin^2\theta)}{M+m}>0.
\]
约化运动方程 \(m_{\rm eff}\ddot s+ks=mg\sin\theta\)，平衡位置 \(s_{\rm eq}=mg\sin\theta/k\)，角频率 \(\omega=\sqrt{k/m_{\rm eff}}\)。初始条件给
\[
s(t)=s_{\rm eq}+(s_0-s_{\rm eq})\cos\omega t,\quad
X(t)=X_0-\frac{m\cos\theta}{M+m}[s(t)-s_0].
\]
楔块向左运动以平衡滑块向右的水平动量。

(3) 因重力势能为 \(-mg s\sin\theta\)，总机械能化简为
\[
E=\frac12m_{\rm eff}\dot s^2+\frac k2(s-s_{\rm eq})^2-\frac k2s_{\rm eq}^2.
\]
对 \(t\) 求导恰为 \(\dot s(m_{\rm eff}\ddot s+k[s-s_{\rm eq}])=0\)。当 \(M\to\infty\)，\(m_{\rm eff}\to m,\ X(t)\to X_0\)，恢复固定斜面上的单自由度弹簧振子；但当 \(M\) 有限，直接把滑块惯性质量写成 \(m\) 会遗漏楔块后坐。

(4) 单位法向量取指向斜面外的 \(\boldsymbol n=(\sin\theta,\cos\theta)\)。滑块加速度是 \((\ddot X+\ddot s\cos\theta,-\ddot s\sin\theta)\)，其法向投影为 \(\ddot X\sin\theta\)。因此
\[
N-mg\cos\theta=m\ddot X\sin\theta,\quad
N=m[g\cos\theta+\ddot X\sin\theta].
\]
又 \(\ddot X=-m\cos\theta\,\ddot s/(M+m)\)、\(\ddot s=-\omega^2(s_0-s_{\rm eq})\cos\omega t\)，故
\[
N(t)=mg\cos\theta+
\frac{m^2\sin\theta\cos\theta}{M+m}\omega^2(s_0-s_{\rm eq})\cos\omega t.
\]
对整个周期都保持接触的必要且充分（理想单侧法向模型）条件是
\[
g\cos\theta\ge\frac{m\sin\theta\cos\theta}{M+m}\omega^2|s_0-s_{\rm eq}|.
\]
若此条件不满足，部分时段出现负的形式法向反力，实际滑块会脱离斜面，约束运动解不再适用。
